\section{PART I --- SMOOTH MANIFOLDS AND RIEMANNIAN GEOMETRY}

\subsection{1. Smooth Manifolds}
\begin{itemize}
\item Charts, atlases, smooth maps; transition functions
\item Tangent spaces and tangent bundle; differential of a map
\item Vector fields; flows; Lie bracket
\item Submanifolds; Whitney embedding theorem; examples (spheres, Lie groups, Grassmannians)
\end{itemize}

\subsection{2. Differential Forms and Stokes' Theorem}
\begin{itemize}
\item Exterior algebra; $k$-forms; wedge product
\item Exterior derivative; Poincar\'e lemma; de Rham cohomology
\item Stokes' theorem (general); Green, Gauss, classical Stokes as special cases
\item Integration over chains; orientation
\end{itemize}

\subsection{3. Vector Bundles and Connections}
\begin{itemize}
\item Vector bundles; trivial and non-trivial; sections
\item Connections; parallel transport; holonomy; geodesics
\item Curvature as obstruction to parallel transport; curvature 2-form
\end{itemize}

\subsection{4. Riemannian Geometry}
\begin{itemize}
\item Riemannian metrics; Levi-Civita connection; geodesic equation
\item Exponential map; normal coordinates; cut locus
\item Curvature: sectional, Ricci, scalar; geometric meaning
\item Comparison theorems: Bonnet-Myers; Rauch; applications
\end{itemize}

\subsection{5. Lie Groups and Symmetric Spaces}
\begin{itemize}
\item Lie groups: $SO(n)$, $SU(n)$, $SE(n)$, $Sp(n)$, $GL(n)$
\item Lie algebra; exponential map; adjoint representation; BCH formula
\item Homogeneous spaces; Riemannian symmetric spaces
\item Applications: statistics on SPD matrices, Grassmannians, robotics
\end{itemize}

\subsection{6. Information Geometry}
\begin{itemize}
\item Statistical models as Riemannian manifolds; Fisher information metric
\item $\alpha$-connections; duality; Amari's framework
\item Dually flat manifolds; exponential and mixture families
\item Natural gradient descent; connections to optimization (Course 5)
\item Cram\'er-Rao bound from the geometric perspective
\end{itemize}

\section{PART II --- ALGEBRAIC TOPOLOGY}

\subsection{7. Homotopy Theory}
\begin{itemize}
\item Homotopy equivalence; fundamental group; van Kampen theorem
\item Covering spaces; deck transformations; universal cover
\item Higher homotopy groups (brief); Whitehead theorem
\item Applications: robotics; configuration spaces; topological phases of matter
\end{itemize}

\subsection{8. Homology and Cohomology}
\begin{itemize}
\item Simplicial complexes; chain complexes; simplicial homology
\item Singular homology; Mayer-Vietoris; long exact sequence
\item De Rham cohomology; de Rham theorem
\item Poincar\'e duality; applications to manifold topology
\end{itemize}

\subsection{9. Characteristic Classes}
\begin{itemize}
\item Vector bundles; classifying spaces; universal bundles
\item Chern classes (complex); Pontryagin classes (real); Whitney sum formula
\item Euler class; Gauss-Bonnet-Chern theorem
\item Applications: topological insulators; anomaly cancellation; index theorem (brief)
\end{itemize}

\section{PART III --- TOPOLOGICAL DATA ANALYSIS}

\subsection{10. Persistent Homology}
\begin{itemize}
\item Filtrations of simplicial complexes; nested inclusions
\item Persistence modules; barcodes and persistence diagrams
\item Stability theorem: bottleneck distance; 1-Wasserstein stability
\item Algorithms: Vietoris-Rips; \v{C}ech; alpha complexes; Ripser, GUDHI
\end{itemize}

\subsection{11. Topological Summaries}
\begin{itemize}
\item Persistence images; persistence landscapes; Betti curves
\item Mapper algorithm: topological summaries of high-dimensional data
\item Euler characteristic transforms; integral geometry
\item Applications: neuroscience; materials science; biology; time series
\end{itemize}

\subsection{12. Sheaves on Networks}
\begin{itemize}
\item Sheaves on simplicial complexes; sheaf cohomology
\item Sheaf Laplacian; cellular sheaves
\item Applications: sensor networks; distributed inference; GNNs
\end{itemize}

\subsection{13. Topology and Machine Learning}
\begin{itemize}
\item Topological loss functions for neural networks
\item Topology of loss landscapes; critical points
\item Manifold hypothesis; intrinsic dimensionality estimation
\item Persistent homology for time series and natural language
\end{itemize}

\section{PART IV --- GEOMETRIC MECHANICS AND DEEP LEARNING}

\subsection{14. Symplectic Geometry}
\begin{itemize}
\item Symplectic manifolds; Darboux theorem
\item Hamiltonian vector fields; Poisson brackets; Liouville's theorem
\item Moment maps; symplectic reduction; Marsden-Weinstein
\item Symplectic integrators (connection to Courses 3 and 4)
\end{itemize}

\subsection{15. Discrete Differential Geometry}
\begin{itemize}
\item Discrete exterior calculus on simplicial meshes
\item Discrete curvature: angle defect; mean curvature via Laplace-Beltrami
\item Cotangent Laplacian; applications to geometry processing
\item Surface PDEs and heat flow on meshes
\end{itemize}

\subsection{16. Geometric Deep Learning}
\begin{itemize}
\item Symmetries as geometric priors; equivariance as a design principle
\item Group actions on function spaces; invariant and equivariant maps
\item Graph neural networks: message passing as discrete exterior calculus
\item Gauge-equivariant networks; principal bundle formalism
\item $SO(3)$ and $SE(3)$ equivariant architectures: spherical CNNs, E3NN, NequIP
\item Connections to representation theory (Course 2, Sections 32--33)
\end{itemize}
